\title{Parameter-Efficient Orthogonal Finetuning via Factorization}

\begin{abstract}
The prevalence of large foundation models has grown significantly, yet their initial training remains highly resource-intensive. Consequently, strategies that can adapt these advanced models to new tasks in an efficient manner have become increasingly crucial. This work explores a robust approach to model adaptation known as Orthogonal Finetuning (OFT). While OFT is recognized for its solid transferability, its reliance on orthogonal matrices leads to a substantial parameter footprint due to their high dimensional nature. To mitigate this, we reinterpret OFT through the lens of information flow, elucidating key criteria essential for enhancing parameter efficiency. Drawing an analogy to the efficient signal processing achieved by the Cooley-Tukey fast Fourier transform algorithm, we introduce a parameterization technique utilizing butterfly structures to realize orthogonal transformations more efficiently. Integrating this factorization into OFT yields Orthogonal Butterfly~(BOFT), a novel finetuning strategy with marked improvements in parameter efficiency. BOFT extends OFT, offering a broader and more generalizable orthogonal finetuning paradigm. We empirically validate the effectiveness of BOFT on a range of downstream tasks in computer vision and natural language processing, including adaptation of large vision transformers, language models, and text-to-image diffusion architectures.
\end{abstract}

\section{Introduction}

Breakthroughs such as ChatGPT~\cite{brown2020language,chen2021evaluating} and Stable Diffusion~\cite{rombach2022high} exemplify the extraordinary generalization performance enabled by large foundation models. The exceptional success of these systems has coincided with explosive growth in model size—GPT-3, for instance, comprises approximately 175 billion parameters—which in turn renders training entirely new foundation models an increasingly daunting endeavor. As such, much of the research community's progress hinges on methods for adapting pretrained foundation models rather than building them anew. To enable the practical transfer of these large models to downstream applications, several adaptation methodologies have been developed: \emph{model finetuning}~\cite{hulora2022,zaken2022bitfit,guo2020parameter,raffel2020exploring,qiu2023controlling,zhang2022adaptive,chavan2023one}, which leaves the core model structure intact while selectively updating a subset of its weights; \emph{adapter tuning}~\cite{rebuffi2017learning,houlsby2019parameter,pfeiffer2020adapterfusion,lin2020exploring,he2021towards}, where extra trainable modules are incorporated alongside the original parameters; and \emph{prompt tuning}~\cite{li2021prefix,lester2021power}, which involves learning additional input-specific tokens as part of the adaptation process. Among these adaptation frameworks, model finetuning stands out for its straightforward implementation, strong empirical results, and crucially, it avoids introducing additional latency during inference.

Model finetuning fundamentally aims to maintain the similarity between the original pretrained model and its finetuned counterpart according to specific criteria, thereby preserving the pre-existing knowledge acquired during pretraining. In practice, most existing finetuning strategies employ a reduced learning rate, effectively limiting the changes made to the model and keeping the deviation between the pretrained and finetuned states minimal. Recognizing the inherent structure within weight matrices, a more systematic method involves maintaining the relational structure within these matrices—specifically, the angles between pairs of neurons. This motivates the design of Orthogonal Finetuning~(OFT)~\cite{qiu2023controlling}, a new approach in which an orthogonal matrix is used to transform all neurons within a single layer identically, thereby mathematically ensuring that the angles between neurons are preserved during finetuning. Although OFT has achieved strong results in terms of both generalization and optimization when applied to text-to-image diffusion model finetuning~\cite{qiu2023controlling}, its practicality is hindered by the substantial parameter count required for large orthogonal matrices. To mitigate this, OFT proposes a block-diagonal formulation, which partitions the orthogonal transformation into independent blocks, reducing parameter overhead. However, this comes with notable drawbacks: the fixed block-diagonal structure imposes a rigid sparsity pattern, applies transformations separately in each block, limits model expressiveness (for instance, making it incapable of approximating certain conventional linear operations), and, crucially, does not provide a principled method for selecting an optimal sparsity structure.

The central challenge lies in constructing a dense orthogonal matrix while maintaining parameter efficiency. Conventionally, a fully dense orthogonal matrix in $d$ dimensions entails $\mathcal{O}(d^2)$ parameters, seemingly making this task prohibitive. However, our approach introduces a novel methodology: we compose a dense orthogonal matrix from a sequence of sparse, factorized matrices. This strategy exploits the notion that a reduction in parameter count can be achieved by substituting additional computational effort for memory usage. Representing the orthogonal transformation as a product of sparse matrices means that, during training, repeated matrix multiplication is necessary in every iteration. To contextualize this factorization, we reinterpret dense orthogonal matrix generation as a form of information transmission. Under this lens, synthesizing a dense orthogonal matrix from sparse factors becomes analogous to communicating information over a grid-like graph structure. Such an information-theoretic viewpoint paves the way for a multitude of possible sparse matrix factorizations that keep parameter numbers low yet retain sufficient expressiveness to represent dense matrices.

For parameter reduction within this information transmission framework, we are inspired by the butterfly network topology found in the Cooley-Tukey fast Fourier transform algorithm~\cite{cooley1965algorithm}, where efficient information propagation among nodes is a key feature~\cite{klasing1994broadcasting}. The butterfly graph intrinsic to Cooley-Tukey provides a natural scheme for sparse matrix decomposition—referred to as butterfly factorization. Utilizing butterfly factorization allows the construction of a dense $d$-dimensional matrix via a product of $\mathcal{O}(\log d)$ sparse factors, with each sparse matrix possessing only $\mathcal{O}(d)$ nonzero entries. By ensuring each of these matrices is orthogonal, we attain an efficient orthogonal parameterization, requiring merely $\mathcal{O}(d\log d)$ parameters—a drastic reduction compared to the full parameter count. Capitalizing on this efficient structure, we introduce a new finetuning paradigm, Orthogonal Butterfly~(BOFT), which serves as a generalized, parameter-efficient orthogonal finetuning method. OFT emerges as a specific instance within the broader BOFT framework. Both the block-diagonal and butterfly structures share the property of inducing sparsity in orthogonal matrices, thereby limiting the number of trainable parameters. It is instructive to compare our framework to the low-rank structure employed in LoRA~\cite{hulora2022}; low-rank and sparse matrices represent two prominent categories of structured matrices~\cite{candes2011robust} aimed at improving parameter efficiency.

In contrast to OFT's use of a block-diagonal configuration to balance model expressivity and structural regularity, BOFT leverages the butterfly pattern to achieve a more continuous interpolation between matrices drawn from the entire orthogonal group (for expressivity) and identity matrices (for regularity). This approach allows for the discovery of a more compact set of hypotheses within the orthogonal group suited to downstream tasks. Given the butterfly structure's prevalence in efficient linear operations—such as those found in discrete Fourier and discrete cosine transforms—we posit that this specific parameter-efficient architecture induces a beneficial implicit inductive bias, thereby enhancing model generalization and mitigating overfitting. Our rationale is rooted in the observation that the butterfly configuration readily encompasses many canonical linear transforms, a capability absent in the block-diagonal construction utilized by OFT. Our principal contributions are outlined below:

\begin{itemize}[leftmargin=*,nosep]
    \item We introduce a new information transmission perspective on parameter-efficient orthogonal finetuning. This reframing converts the challenge of designing a dense, parameter-efficient orthogonal matrix into analyzing information flow across a grid-based graph structure.
    \item Drawing from the butterfly architecture inherent in the Cooley-Tukey algorithm, we put forth Orthogonal Butterfly, an orthogonal finetuning technique that achieves parameter efficiency within the confines of the information transmission framework.
    \item We provide theoretical explanations illustrating how BOFT markedly decreases the number of trainable parameters, yet preserves robust expressivity and stable optimization during training. Furthermore, matrix factorization in BOFT gives rise to an appealing mechanism for weight interpolation (see Figure~\ref{fig:boft_interpolation}).
    \item For the first time, we extend the application of orthogonal finetuning to diverse domains beyond controllable text-to-image generation~\cite{qiu2023controlling}, thereby demonstrating its broad potential as a general-purpose model finetuning paradigm. To substantiate this, we evaluate BOFT on a spectrum of adaptation tasks in both computer vision and natural language processing, where it substantially surpasses leading parameter-efficient methods, underscoring its enhanced efficiency and generalization performance.
\end{itemize}

\textbf{Parameter-efficient finetuning (PEFT)}. The escalating size and capabilities of foundation models (\emph{e.g.}, \cite{brown2020language,radford2021learning,rombach2022high,kirillov2023segment}) have led to a growing body of research focused on developing parameter-efficient strategies for their finetuning on downstream applications~\cite{treviso2023efficient,ding2023parameter,lialin2023scaling}. Numerous approaches for PEFT have emerged~\cite{houlsby2019parameter,aghajanyan2020intrinsic,hulora2022,edalati2022krona,wang2022adamix,gheini2021cross,zaken2022bitfit,guo2020parameter,sung2021training,ansell2022composable,lester2021power,li2021prefix,vu2022spot,he2021towards,mao2021unipelt,karimi2021compacter,liu2022few,sung2022lst,chen2022parameter,jia2022visual,chen2022adaptformer,zhang2022neural,jie2023fact,lian2022scaling,luo2023towards}, with reparameterization-based methods~\cite{aghajanyan2020intrinsic,hulora2022,edalati2022krona} being particularly pertinent to this paper. For example, LoRA~\cite{hulora2022} introduces updates to the pretrained weight matrix via the addition of a product between two low-rank matrices, resulting in effective performance on a range of natural language processing tasks. In the original LoRA, the added matrices’ rank remains fixed, but several later works~\cite{zhang2022adaptive,valipour2022dylora,zhang2023increlora} introduce methods to dynamically tune the rank for each layer, thereby optimizing the use of available parameters. The widespread adoption of such low-rank weight reparameterizations can be attributed to their straightforward implementation~\cite{dettmers2023qlora,zi2023delta,chavan2023one}. Building on insights regarding the connection between hyperspherical energy and model generalization~\cite{liu2018learning,liu2021orthogonal}, orthogonal finetuning has been introduced in~\cite{qiu2023controlling} as an alternative approach for adapting text-to-image diffusion models, where an orthogonal matrix is learned to rotate the neurons within the same layer. This method yields both improved generalization and greater training stability compared to LoRA. However, OFT typically involves training more parameters than LoRA, motivating further work to improve its parameter efficiency. Additionally, it remains an open question whether OFT can generalize effectively to adaptation tasks beyond those considered in~\cite{qiu2023controlling}. BOFT advances this line of work by leveraging butterfly factorization to enhance OFT’s parameter efficiency. This development enables a systematic exploration of orthogonal finetuning in a broader context of model adaptation tasks.

\textbf{Butterfly structures}. The radix-2 Cooley-Tukey algorithm efficiently computes the $N$-point discrete Fourier transform by recursively decomposing it into two smaller transforms of size $\frac{N}{2}$, leading to a “butterfly” computational pattern characterized by products of several sparse matrices, collectively termed a butterfly matrix. These butterfly matrices have previously been employed in various contexts, including providing an efficient parameterization of orthogonal matrices to circumvent pivoting in Gaussian elimination~\cite{parker1995random}, supporting the stability of recurrent neural network training~\cite{jing2017tunable}, and aiding in kernel approximations~\cite{munkhoeva2018quadrature}. Furthermore, butterfly parameterizations facilitate fast linear transformations~\cite{dao2019learning,dao2019kaleidoscope}, and have been utilized to accelerate the training of neural networks~\cite{chen2022pixelated,dao2022monarch}. Beyond these, butterfly structures play important roles in speeding up matrix-vector products~\cite{michielssen1996multilevel,o2010algorithm}, compressing data-sparse matrices~\cite{li2015butterfly}, and improving network data transmission~\cite{klasing1994broadcasting,li2003linear,rajasingh2016transmission}. In this study, unlike prior works, our focus is on deploying butterfly structures to make OFT more parameter-efficient.

\section{An Information Transmission View on Orthogonal Finetuning}

We begin by outlining the foundational concepts of OFT. In order to adapt pretrained weights, OFT introduces a reparameterization scheme where the new weight matrix is expressed as the multiplication of a trainable orthogonal matrix with the original fixed weight matrix. Unlike LoRA, which implements additive modifications via a low-rank matrix, OFT updates weights using a multiplicative orthogonal transformation. LoRA achieves parameter efficiency through its low-rank decomposition, whereas standard OFT employs an orthogonal matrix with a block-diagonal form~\cite{qiu2023controlling}; reducing the block size increases OFT's parameter efficiency. Figure~\ref{fig:comp_lora} visually contrasts these approaches. The rationale behind utilizing an orthogonal transformation for fine-tuning lies in maintaining the pairwise angular relationships among neuron vectors~\cite{liu2018learning,liu2021orthogonal,qiu2023controlling}, thus safeguarding much of the semantic information acquired during pretraining.

More specifically, OFT learns an orthogonal matrix $\bm{R} \in \mathbb{R}^{d\times d}$ for a pretrained linear transformation $\bm{W}^0\in\mathbb{R}^{d\times n}$. This alters the forward computation from $\bm{z} = (\bm{W}^0)^\top\bm{x}$ to $\bm{z} = (\bm{R}\bm{W}^0)^\top\bm{x}$, where $\bm{x}\in\mathbb{R}^{d}$ represents the input and $\bm{z}\in\mathbb{R}^{n}$ denotes the output. Zero initialization of OFT is realized by setting $\bm{R}$ initially to the identity matrix. To guarantee $\bm{R}$ remains orthogonal throughout training, we adopt the Cayley parameterization, following \cite{qiu2023controlling,liu2021orthogonal}, so that $\bm{R} = (\bm{I} + \bm{Q})(\bm{I} - \bm{Q})^{-1}$ where $\bm{Q}$ is skew-symmetric, satisfying $\bm{Q} = -\bm{Q}^\top$. To enhance parameter efficiency, the orthogonal matrix $\bm{R}$ is structured as a block-diagonal matrix, written as $\text{diag}(\bm{R}_1, \bm{R}_2, \cdots, \bm{R}_r)$, where each $\bm{R}_i\in\mathcal{R}^{b\times b}$ is an independent small orthogonal matrix, and $br=d$.

However, the adoption of the block-diagonal form introduces an implicit premise: each neuron's dimension (a column in $\bm{W}^0$) is partitioned into $r$ separate groups, each processed with its own orthogonal transformation. Although empirical results demonstrate the effectiveness of the block-diagonal configuration, assigning neuron dimensions to different groups solely by their indices lacks theoretical justification, which underscores the potential advantages of using dense orthogonal matrices. This leads to a fundamental inquiry: \emph{Is it feasible to design a dense orthogonal matrix that retains parameter-efficiency?}

In response to this problem, we introduce a method that constructs a dense orthogonal matrix by composing it as a sequence of several sparse orthogonal matrices. To unify and clarify the study of sparsity in orthogonal matrix decompositions, we recast the task of synthesizing a dense orthogonal matrix in OFT as analogous to an information transmission scenario. More precisely, building a dense matrix $\bm{R}\in\mathbb{R}^{d\times d}$ through the multiplication of $m$ square matrices, expressed as $\thickmuskip=2mu \medmuskip=2mu \bm{R}=\bm{B}_m\bm{B}_{m-1}\cdots\bm{B}_1$, is conceptualized as routing information across a $d\times (m+1)$ grid network, depicted in Figure~\ref{fig:info_trans}. This transmission perspective arises from viewing a $d$-dimensional dense square matrix as representing complete connections from one group of $d$ nodes to another. In this analogy, if $R_{ij}$ in $\bm{R}$ is zero, it signifies the absence of a communication channel from node $j$ to node $i$, whereas a non-zero entry implies connectivity. Thus, decomposing the dense matrix $\bm{R}$ into products of sparse factors $\bm{B}_m\bm{B}_{m-1}\cdots\bm{B}_1$ equates to stepwise information exchange structured by the respective graphs of $\bm{B}_i$ for each $i$. Information transfer starts with the structure imposed by $\bm{B}_1$ and proceeds sequentially to $\bm{B}_n$ as the final stage. As detailed in Figure~\ref{fig:info_trans}, the example of $\bm{R}=\bm{B}_5\bm{B}_4\bm{B}_3\bm{B}_2\bm{B}_1$ is presented, with each factor’s sparsity layout and its corresponding graph visualization. The illustrated graph results from expanding the multiplication process, where $\bm{B}_i$ serves as the adjacency matrix connecting the $i$-th layer of nodes to those in the $(i+1)$-th layer. Explicitly, the entry at position $(j_1,j_2)$ in $\bm{B}_i$ indicates the presence (or absence, if zero) of a directed link from node $j_2$ at level $i$ to node $j_1$ at level $(i+1)$. Achieving a fully populated $\bm{R}$ through the product $\bm{B}_5\bm{B}_4\bm{B}_3\bm{B}_2\bm{B}_1$ requires that every node at the initial level must be capable of relaying information to all nodes at the sixth level. For the truncated sequence $\bm{R}=\bm{B}_4\bm{B}_3\bm{B}_2\bm{B}_1$, corresponding to nodes spanning from the first to fifth levels, one observes, for instance, that node 1 cannot transmit to node 3, reflecting a zero in position $(3,1)$ in the combined sparsity pattern. This pattern of correspondence between graph connectivity and sparsity is consistent in shorter products, such as $\bm{R}=\bm{B}_3\bm{B}_2\bm{B}_1$ and $\bm{R}=\bm{B}_2\bm{B}_1$. More generally, for a matrix $\bm{R}\in\mathbb{R}^{d\times d}$ to attain full density, its factorization via matrices $\bm{B}_m,\cdots,\bm{B}_1$ must generate a network of directed edges across a $d\times (m+1)$ lattice. Here, each edge connects only two nodes at consecutive levels (that is, columns), and a viable transmission pathway from each node in the first level to every node in the $(m+1)$-th level must exist.

The information transmission perspective offers valuable insights into the sparse structure present in factorization matrices within OFT. In Figure~\ref{fig:blk_diag}, the block-diagonal arrangement of $\bm{R}$, characteristic of standard OFT, is illustrated. While this organization decreases the count of trainable parameters, it inherently lacks the capacity to generate a fully dense matrix $\bm{R}$. Our objective is to assemble a dense orthogonal matrix by combining $m$ sparse orthogonal matrices, all while minimizing the count of effective trainable parameters required. According to the information transmission viewpoint, our main criteria for this construction are: \emph{(i)} \textbf{dense connectivity}, meaning every node from the input layer must be linked to every output node through at least one pathway; and \emph{(ii)} \textbf{minimal free edges}, where the overall quantity of edges should be as limited as possible without breaching orthogonality. Importantly, orthogonality imposes specific restrictions on the inter-level connections. To ensure each $\bm{B}_i$ maintains full rank—a prerequisite for orthogonality—it is necessary to connect $d$ edges between the $i$-th and $(i=1)$-th levels, establishing a one-to-one correspondence among all respective nodes. This requirement dictates that a minimum of $d$ edges must connect adjacent levels (as seen in Figure~\ref{fig:info_trans} for the case $d=4$). Since these $d$ edges are essential for preserving orthogonality and are not learnable (e.g., in a $d \times d$ orthogonal matrix, $d$ non-zero entries are constrained to $\pm 1$), they are excluded from the count of trainable edges. The orthogonality constraint thus introduces dependencies among connections, resulting in a reduced parameter count compared to a general dense matrix. For instance, in a $2\times 2$ orthogonal matrix, setting the $(1,1)$ entry to zero mandates that the $(2,2)$ entry must also be zero, meaning the removal of one edge may require the elimination of another. Consequently, the admissibility of edge configurations can shift as orthogonality may force certain edges to be present or absent. Visualizing the sparsity pattern of the resulting orthogonal matrix, this information transmission approach is especially suited to OFT, as our concern is only with the trainable non-zero components of $\bm{R}$, and not their exact numerical values.

A straightforward dense connection between two hierarchical levels would necessitate $\mathcal{O}(d^2)$ edges, that is, a fully connected orthogonal matrix comprising $d^2-d$ off-diagonal edges (for instance, $d=4$ leads to 12 edges). In contrast, the factorization approach depicted in Figure~\ref{fig:info_trans} employs only $10$ edges overall, which is fewer than required for a single dense orthogonal matrix. This conceptual framework facilitates the exploration of the parameter efficiency of OFT through a graph-based lens, readily allowing one to devise valid matrix factorizations. Inspired by the characteristic connectivity of butterfly graphs~\cite{cooley1965algorithm} from the Cooley-Tukey algorithm, this topology offers an efficient scheme, achieving complete bipartite connectivity between $d$ source and $d$ target nodes with just $\mathcal{O}(d\log d)$ edges. For illustration, the structure shown in Figure~\ref{fig:info_trans} provides full connectivity with $10$ edges, while the butterfly network topology achieves the same goal with only $8$ edges. We proceed to detail how employing the butterfly structure enhances parameter efficiency.

\section{Orthogonal Parameterization by Butterfly Factorization}

The butterfly structure, central to the Cooley-Tukey algorithm, underpins the fast Fourier transform. In the context of Fourier analysis, alterations in one domain (frequency) can result in widespread modifications across another domain (space), paralleling the information transmission context here—where each node at the initial level has pathways to all nodes in the final level. Beyond signal processing, the butterfly structure is prominent as a network topology in computer systems~\cite{solihin2015fundamentals,li2011linear} due to its efficacy in facilitating data exchange. Letting $k \geq 2$ denote a power of $2$, we define the \emph{butterfly factor} $\bm{B}^F(k)$ as follows:

\begin{equation}\label{eq:but_factor}
\begin{aligned}
    \bm{B}^F(k)=\begin{bmatrix}
    \text{diag}(\bm{d}_1) & \text{diag}(\bm{d}_2) \\
    \text{diag}(\bm{d}_3) & \text{diag}(\bm{d}_4)
    \end{bmatrix}\in\mathbb{R}^{k\times k},
\end{aligned}
\end{equation}

where each $\bm{d}_i\in\mathbb{R}^{\frac{k}{2}}$ for $i=1,\ldots,4$ represents a vector. For $d=2^N$, we recursively construct the $d$-dimensional \emph{butterfly matrix} $\bm{B}(d)\in\mathbb{R}^{d\times d}$ as:

\begin{equation}
\begin{aligned}
    \!\!\bm{B}(d)=\tilde{\bm{B}}(d,d)\!\cdot\!\begin{bmatrix}
    \bm{B}_1(\frac{d}{2})\!\!\!\!\! & \bm{0} \\
    \bm{0}\!\!\!\!\! &\bm{B}_2(\frac{d}{2})
    \end{bmatrix}= \tilde{\bm{B}}(d,d)\tilde{\bm{B}}(d,\frac{d}{2})\cdots\tilde{\bm{B}}(d,2),
\end{aligned}
\end{equation}

where $\bm{B}_1(\frac{d}{2})$ and $\bm{B}_2(\frac{d}{2})$ denote two butterfly matrices, each operating in $\frac{d}{2}$ dimensions. We introduce the \emph{butterfly component} as $\thickmuskip=2mu \medmuskip=2mu \tilde{\bm{B}}(d,k)=\text{diag}(\bm{B}^F_1(k),\cdots,\bm{B}^F_{d/k}(k))$, which forms a block-diagonal matrix of dimensions $d \times d$, with each block of size $k$. Here, each $\bm{B}^F_i(k)$, for all $i$, refers to a butterfly factor as given in Equation~\ref{eq:but_factor}. This framework enables the use of butterfly matrices to parameterize an orthogonal matrix, provided all constituent factors $\tilde{\bm{B}}(d,k)$ within the overall butterfly matrix $\bm{B}(d)$ maintain orthogonality. Considering specifically the block-diagonal matrix $\tilde{\bm{B}}(d,2)$ (with $2\times 2$ blocks), orthogonality can be assured by employing the Cayley transform or by parameterizing each block via $2$-dimensional rotations, following techniques in \cite{liu2021orthogonal,qiu2023controlling}. Since the non-zero structure of any butterfly component represents a permutation of that in $\tilde{\bm{B}}(d,2)$, this parameterization method extends straightforwardly to all such components, resulting in an efficient approach to construct orthogonal matrices from multiple $2 \times 2$ orthogonal blocks. To further generalize, adopting the approach from \cite{chen2022pixelated}, we define a block butterfly component $\tilde{\bm{B}}^b(d,k)$ in which every $\bm{d}_i$ entry is substituted with a $b \times b$ matrix. Ensuring that the block butterfly component $\tilde{\bm{B}}^b(d,2)$ is orthogonal involves parameterizing each $2b\times 2b$ block as an orthogonal matrix. For components $\tilde{\bm{B}}^b(d,k)$ with $k > 2$, their block-wise sparsity pattern arises as a permutation of that in $\tilde{\bm{B}}^b(d,2)$, so the same orthogonality-preserving strategy applies. Integrating these insights, the BOFT forward computation becomes

\begin{equation}
    \begin{aligned}\nonumber
    \bm{z} = \big(\bm{R}(m,b)\cdot\bm{W}^0\big)^\top\bm{x},~~~~\text{s.t.}~\bigg{\{}\bm{R}(m,b) = \prod_{i=1}^m \tilde{\bm{B}}^b_{(d,i)}~~\&~~\underbrace{\big(\tilde{\bm{B}}^b_{(d,j)}\big)^\top\tilde{\bm{B}}^b_{(d,j)}=\tilde{\bm{B}}^b_{(d,j)}\big(\tilde{\bm{B}}^b_{(d,j)}\big)^\top\!=\bm{I}_d}_{\forall j\in [1, m]}\bigg{\}},
    \end{aligned}
\end{equation}

For notational simplicity, we refer to $\tilde{\bm{B}}^b(d,2^{m-i+1})$ as $\tilde{\bm{B}}^b_{(d,i)}$, with $\bm{I}_d$ indicating the $d\times d$ identity matrix. The orthogonal matrix $\bm{R}(m,b)\in\mathbb{R}^{d\times d}$ is constructed by multiplying multiple orthogonal butterfly factors. For clarity, BOFT using $\bm{R}(m,\frac{b}{2})$ is denoted as BOFT($m$,$b$) for $b\geq 2$. When setting $m=1$, BOFT($1$,$b$) reduces to the block-diagonal OFT~\cite{qiu2023controlling} with block size $b$. If we further take $b=d$, BOFT($1$,$d$) collapses to the conventional OFT~\cite{qiu2023controlling} in which $\bm{R}$ is a full orthogonal matrix. Selecting $m=\log \frac{2d}{b}$ with block size $b$ causes BOFT to be equivalent to using a dense block butterfly matrix $\bm{B}^{\frac{b}{2}}(d)$ as $\bm{R}$. Typically, BOFT($m$,$b$) entails $\frac{1}{2}(b-1)dm$ learnable parameters, corresponding to the case where one finetunes a linear layer of dimensions $d\times n$. For the butterfly matrix setup ($m=\log d,\,b=2$), the parameter cost for BOFT is $\mathcal{O}(d\log d)$. This contrasts with the original OFT approach, which demands $\mathcal{O}(d^2)$ parameters for a dense orthogonal matrix, and the block-diagonal OFT, which has a parameter count of $\mathcal{O}(bd)$ for $r$ blocks. As a result, constructing a dense orthogonal matrix in the original OFT requires setting $b=d$, while BOFT allows one to attain such matrices for any $b$.

\textbf{Identity initialization for BOFT}. In standard finetuning practice, one often initializes from the pretrained model itself to ensure the adaptation remains close to the original. For instance, LoRA initializes the added low-rank weights with zeros. Analogously, in BOFT, each butterfly component is set as the identity matrix during initialization (in the Cayley transform, this amounts to initializing skew-symmetric matrices as zero).

\textbf{Multiplicative Dropout for BOFT}. LoRA~\cite{hulora2022} applies Dropout to the low-rank adaptation modules to combat overfitting. The typical Dropout method~\cite{srivastava2014dropout} integrates naturally with LoRA, but is not directly suitable for BOFT given the multiplicative update structure of the weights. To tackle this, we introduce a multiplicative Dropout mechanism tailored to BOFT. Given that $\bm{R}(m,b)$ is factored into $m$ orthogonal butterfly elements that can be reordered into $2b\times 2b$ block-diagonal orthogonal matrices, multiplicative Dropout operates by selecting a random $p_1$ fraction of butterfly components and a $p_2$ fraction of diagonal blocks within each component. These selected blocks are then replaced by identity matrices, thereby encouraging regularization.

\section{Intriguing Insights and Discussions}

\textbf{Expressivity of BOFT}. The combination of the butterfly structure and permutations enables the recovery of several well-known fast linear transforms~\cite{dao2019learning,dao2019kaleidoscope} (\emph{e.g.}, fast Fourier transform, Hadamard transform) in an exact manner. However, it remains unclear how effectively our orthogonal butterfly matrix approximates arbitrary orthogonal matrices. To address this, we conducted an experiment aimed at approximating a randomly generated dense orthogonal matrix~\cite{anderson1987generation} of size $512\times 512$. Results, summarized in Figure~\ref{fig:para_eff}, are averaged across 10 random initializations. The y-axis presents the approximation error while the x-axis corresponds to the count of effective, trainable parameters. Each color-coded curve represents BOFT with a particular block size; notably, the leftmost point depicts the error associated with BOFT($1$,$b$) (i.e., the vanilla block-diagonal OFT for block size $b$). In comparison to OFT, BOFT generally achieves improved parameter efficiency. For instance, BOFT($9$,$2$) demonstrates greater expressiveness than BOFT($1$,$16$) despite requiring significantly fewer parameters. Configurations of BOFT characterized by smaller $b$ and higher $m$ tend to be even more parameter-efficient. For example, BOFT($6$,$4$) achieves comparable approximation errors to BOFT($2$,$16$) with a reduced parameter count. Overall, the butterfly matrix spans a more constrained yet structured subspace within the orthogonal group (relative to the block-diagonal configuration), making BOFT strictly more expressive than OFT at equivalent block sizes.

\begin{theorem}[Expressivity of BOFT]\label{thm:exp_boft}
    BOFT surpasses OFT in expressivity for the same block size. To approximate any orthogonal matrix of size $d$, a product of butterfly matrices such as $\bm{B}_{d-1,1}(d)\bm{B}^\top_{d-1,2}(d)\cdots\bm{B}_{1,1}(d)\bm{B}^\top_{1,2}(d)$ can be employed, where every $\bm{B}_{i,j}(d)$ is a butterfly matrix.
\end{theorem}

Theorem~\ref{thm:exp_boft} paves the way for a generalized construction in BOFT—specifically, the orthogonal matrix may be defined as $\thickmuskip=2mu \medmuskip=2mu \bm{R}^G(m_1,b_1,m_2,b_2,l)=\bm{R}_{l,1}(m_1,b_1)\bm{R}_{l,2}^T(m_2,b_2)\cdots\bm{R}_{1,1}(m_1,b_1)\bm{R}_{1,2}^T(m_2,b_2)$, where $\bm{R}_{i,j}^T(m,b)$ stands for the orthogonal matrix applied in BOFT. When $m_1=m_2=\log d$, $b_1=b_2=2$, and $l=d-1$, this structure $\bm{R}^G(m_1,b_1,m_2,b_2,l)$ is capable of spanning the full orthogonal group, coinciding with what is known as the kaleidoscope hierarchy~\cite{dao2019kaleidoscope}. It is important to observe, nonetheless, that increased expressivity does not uniformly correspond to enhanced performance during finetuning; for example, full finetuning—despite being universally expressive—frequently underperforms. Consequently, the balance between expressivity and regularization is central for model generalization in finetuning scenarios. By incorporating structural priors, BOFT broadens the parameter space for finetuning, thereby supporting a more optimal compromise between expressiveness and regularization.

\textbf{Spectral properties}. Orthogonal finetuning tends to achieve superior spectral characteristics compared to LoRA, primarily because it maintains the spectral norm of the original weight matrix $\bm{W}^0$ without alteration. To illustrate, consider the singular value decomposition: $\bm{W}^0=\bm{U}\bm{\Sigma}\bm{V}^\top$, where $\bm{U}$ and $\bm{V}$ are orthogonal, and $\bm{\Sigma}$ is diagonal with the singular values. In both OFT and BOFT, an orthogonal transformation $\bm{R}$ is applied to the left, yielding updated weights $\bm{R}\bm{U}\bm{\Sigma}\bm{V}^\top$. This operation leaves the top singular value—and thus the spectral norm of $\bm{W}^0$—unchanged. Preserving this property has been found to significantly enhance both training stability and generalization~\cite{miyato2018spectral,yoshida2017spectral}. Further mathematical analysis of these properties is provided in Appendix~\ref{app:sn_property}.

\textbf{Orthogonal finetuning as learning bilinear similarity}. BOFT can also be interpreted as optimizing a bilinear similarity function of the form $\bm{w}^0_i\bm{R}\bm{x}$, where $\bm{w}^0_i$ denotes the $i$-th neuron (that is, the $i$-th column of $\bm{W}^0$). In this perspective, BOFT is essentially learning the bilinear similarity matrix $\bm{R}$, subject to an orthogonality constraint, which directly relates to the domains of distance metric learning~\cite{xing2002distance} and the theory of bilinear forms~\cite{roman2005advanced}.

\textbf{Inductive bias and generalization in BOFT}. In the context of BOFT, $\bm{R}(m,b)$ typically describes a structured subset within the orthogonal group, which consequently restricts the hypothesis space. This structural constraint imposes an inductive bias on the learning process. We propose that the type of inductive bias generated by butterfly factorization is advantageous for generalization, since it shares structural motifs with a variety of classical linear operators~\cite{dao2019kaleidoscope} such as the discrete Fourier transform, discrete sine and cosine transforms, as well as the Hadamard transform. Additionally, the utilization of sparse matrix factorization within BOFT may introduce further implicit inductive biases~\cite{gunasekar2017implicit,li2018algorithmic,liu2019neural}.

\textbf{Comparison to butterfly-based sparse training}. Several studies~\cite{dao2019learning,dao2019kaleidoscope,chen2022pixelated,dao2022monarch} have explored sparse training via butterfly parameterizations, generally focusing on directly reparameterizing neural network weight matrices as butterfly structures and training these networks from the beginning. In particular, \cite{dao2022monarch} investigates the process of fine-tuning pretrained weights by projecting them onto a modified form of butterfly matrices, subsequently optimizing these projected parameters for specific downstream applications. In contrast, BOFT introduces a fundamentally different fine-tuning approach by applying transformations using layer-shared weight matrices rather than matrix projection.

\input{rebuttal-results}

\section{Concluding Remarks and Limitations}

This work introduces Orthogonal Butterfly, a broadly applicable and parameter-efficient fine-tuning technique for foundation models, utilizing the butterfly matrix structure. The principal advance for achieving higher parameter efficiency lies in representing a dense orthogonal matrix as the product of several sparse orthogonal matrices. To systematically identify suitable matrix decompositions, we develop a graphical information transmission framework. Within this framework, the butterfly configuration emerges as an effective solution for sparse orthogonal matrix factorization. We provide empirical evidence supporting BOFT’s effectiveness across a spectrum of large models, including language models, vision models, and text-to-image generation systems. The experiments consistently indicate that BOFT outperforms other general-purpose fine-tuning methods.

However, BOFT also presents certain limitations. Notably, because the resulting orthogonal matrix is constructed as a product of multiple orthogonal matrices, the training computational cost is marginally higher relative to OFT. Enhancing the training efficiency of BOFT thus remains an open area for future work. On the positive side, once fine-tuning is complete, the orthogonal matrices derived by BOFT can be merged into the pretrained model parameters, ensuring that no additional latency is introduced during inference. Furthermore, it remains an open question whether the butterfly network constitutes the most optimal architecture for information transmission. Our information transmission framework opens up avenues to draw from related disciplines such as computer networking, where analysis of network topology efficiency is well-established. It is anticipated that alternative network structures might further improve the construction of dense orthogonal matrices.